\chapter{Энергия вакуума и горизонты де Ситтера}

\lectureoverview{Энергия вакуума не похожа на обычную материю: она не выделяет предпочтённую инерциальную систему отсчёта и не разрежается по мере расширения Вселенной. Как только учитывается гравитация, этот на первый взгляд безобидный выбор нуля энергии изменяет геометрию пространства-времени. Мы выведем закон при доминировании материи \(a(t)\propto t^{2/3}\), сопоставим его с законом при доминировании вакуума \(a(t)\propto e^{Ht}\), а затем перепишем пространство-время де Ситтера в конформных координатах. Получающаяся причинная диаграмма обнаруживает зависящий от наблюдателя горизонт событий на фиксированном собственном расстоянии.}

\section{Квантовые флуктуации в вакууме}

Начнём с вопроса, который задаёт мотив всей лекции. Квантовая теория поля приписывает пустому пространству энергию. Означает ли это, что реальные электрон--позитронные пары непрерывно возникают из ничего?

\begin{classroomqa}
\audiencequestion{Происходят ли в вакууме в действительности рождение и аннигиляция электрон--позитронных пар, или это лишь удобный способ учёта?\lecturetimestamp{00:00:06}}
\lecturerresponse{Виртуальные пары вносят вклад в энергию вакуума, но это не обычные частицы, за появлением которых можно наблюдать в пустом пространстве. Они представляют собой флуктуации квантового основного состояния и существуют лишь в пределах, допустимых для виртуального процесса.}
\end{classroomqa}

Наиболее наглядную модель даёт гармонический осциллятор. Классически он может оставаться неподвижным в минимуме своего потенциала. Квантовомеханически точное положение в минимуме и точно нулевой импульс несовместимы. Основное состояние является компромиссом между этими двумя требованиями и сохраняет энергию нулевых колебаний
\begin{equation}
E_0=\frac{1}{2}\hbar\omega.
\label{eq:l08-zero-point}
\end{equation}
Это наименьшая энергия, которой может обладать осциллятор. Соответствующее состояние не является возбуждённым и не представляет собой поток наблюдаемых частиц.

Свободное поле разлагается на нормальные моды, каждая из которых ведёт себя как осциллятор. Электромагнитное поле в резонаторе — привычный пример. Каждая мода даёт вклад вида, указанного в уравнении~\eqref{eq:l08-zero-point}; электронное, кварковое и другие квантовые поля вносят собственные флуктуации основного состояния. Суммирование этих вкладов даёт то, что мы называем энергией вакуума. Печально известное несоответствие между наивными микроскопическими оценками и наблюдаемым космологическим масштабом реально, но оно не является темой этой лекции.

\begin{classroomqa}
\audiencequestion{Будет ли ещё обсуждаться расхождение примерно в \(120\) порядков величины?\lecturetimestamp{00:02:59}}
\lecturerresponse{Здесь — нет. Нынешняя цель состоит в том, чтобы понять космологические горизонты: как космологические уравнения приводят к ускоренному расширению и почему ускоренное расширение создаёт горизонты событий.}
\end{classroomqa}

\begin{classroomqa}
\audiencequestion{Является ли ультрафиолет--инфракрасное соотношение из прошлой лекции частью сегодняшнего обсуждения?\lecturetimestamp{00:03:59}}
\lecturerresponse{Возможно, мы вернёмся к нему позже, но в намеченную аргументацию оно не входит. Непосредственный путь ведёт от энергии вакуума к космологии, а затем к горизонтам.}
\end{classroomqa}

\section{Почему энергия вакуума инвариантна относительно преобразований Лоренца}

Следующий вопрос кажется элементарным, но выявляет определяющее свойство энергии вакуума. Если пустое пространство имеет положительную плотность энергии, почему оно не ведёт себя как оптическая среда и не изменяет скорость света?

\begin{classroomqa}
\audiencequestion{Не должна ли положительная плотность энергии вакуума оказывать преломляющее действие, заставляя свет распространяться немного медленнее?\lecturetimestamp{00:04:35}}
\lecturerresponse{Обычное вещество может делать это, потому что задаёт предпочтённую систему покоя. Энергия вакуума иная: она не выделяет такой системы, и каждый инерциальный наблюдатель получает одно и то же значение плотности энергии вакуума. Поэтому локальные измерения света сохраняют инвариантную скорость \(c\).}
\end{classroomqa}

Нужно различать симметрию законов и симметрию конкретного состояния. Уравнения физики могут быть инвариантны относительно трансляций, хотя Stanford — не Berkeley, а Земля — не межзвёздное пространство. Трансляция \emph{всей} физической конфигурации оставляет законы неизменными; реальное распределение материи не обязано выглядеть одинаково в каждой точке. Точно так же уравнения, инвариантные относительно вращений, не делают эту аудиторию вращательно симметричной.

То же различие относится и к преобразованиям Лоренца. У газа есть система покоя. Наблюдатель, движущийся относительно него, видит движущиеся молекулы и другую плотность энергии, так что состояние газа не лоренц-инвариантно, хотя управляющие им уравнения таковы. Обычное вещество, ванны излучения и лабораторная аппаратура — всё это нарушает симметрию, выделяя особые состояния движения или направления.

Энергия вакуума — особая форма материи, которая этого не делает. Представьте себе детектор, показывающий плотность энергии вакуума. В состоянии покоя он даёт некоторое значение \(\rho_{\rm vac}\). Разгоните детектор и повторите измерение: результат не изменится. В релятивистской записи это свойство выражается тензором энергии-импульса, пропорциональным метрике,
\begin{equation}
T_{\mu\nu}^{\rm vac}=-\rho_{\rm vac}g_{\mu\nu},
\qquad
p_{\rm vac}=-\rho_{\rm vac}.
\label{eq:l08-vacuum-stress}
\end{equation}
\editorialnote{Соотношение для давления на доске не выводится. Это ковариантная формулировка двух определяющих свойств, на которых строится лекция: лоренц-инвариантности и постоянной плотности при расширении.}

Локальная оговорка существенна. Для светового луча, проходящего мимо наблюдателя, измерения с помощью близлежащих часов и линеек дают скорость \(c\). Если же спрашивать о скорости далёкого светового луча относительно нас в расширяющейся геометрии, нужно выбрать координаты и задать правило сравнения удалённых событий. Эта координатная скорость не обязана равняться \(c\), хотя каждый локальный наблюдатель измеряет \(c\).

\begin{classroomqa}
\audiencequestion{В этом обсуждении энергия вакуума и космологическая постоянная — синонимы?\lecturetimestamp{00:11:57}}
\lecturerresponse{Да. Это два способа поместить один и тот же член в уравнения Эйнштейна: его можно описывать как космологическую постоянную в геометрической части или как лоренц-инвариантный вакуумный тензор энергии-импульса в части источников.}
\end{classroomqa}

\section{Что разрежается, а что нет}

Заполним большую область нерелятивистскими частицами массы \(m\). Предположим, что на масштабах, достаточно больших для усреднения по галактикам и скоплениям, распределение однородно. Если сопутствующий объём содержит \(N\) частиц, то плотность материи в нём равна
\begin{equation}
\rho_m=\frac{Nm}{V}.
\end{equation}
Когда каждый линейный размер удваивается, объём возрастает в \(2^3=8\), тогда как \(N\) и \(m\) остаются неизменными. Поэтому плотность падает до одной восьмой прежнего значения. Обычная материя разрежается, потому что те же самые частицы занимают больший физический объём.

Излучение тоже разрежается, причём быстрее. Его концентрация падает вместе с объёмом, а длина волны каждого фотона растягивается, и его энергия уменьшается из-за красного смещения. Позже мы запишем результат \(\rho_r\propto a^{-4}\), но решать космологию, в которой доминирует излучение, сегодня вечером не будем.

Энергия вакуума — исключение:
\begin{equation}
\rho_{\rm vac}(t)=\text{constant}.
\label{eq:l08-rho-vac-constant}
\end{equation}
Она не образует сгустков, не выделяет ни направления, ни инерциальной системы отсчёта и не становится менее плотной по мере расширения пространства. Если сопутствующий объём растёт, содержащаяся в нём энергия вакуума растёт пропорционально его объёму.

\editorialnote{Постоянная плотность не означает постоянную полную энергию вакуума внутри расширяющегося сопутствующего объёма. Ковариантный закон сохранения энергии-импульса,
\(\dot\rho+3H(\rho+p)=0\), выполняется потому, что энергия вакуума имеет отрицательное давление \(p_{\rm vac}=-\rho_{\rm vac}\). Член, соответствующий работе, в точности обеспечивает изменение полной энергии вакуума в этом объёме.}

\section{Гравитация чувствительна к выбору нуля энергии}

В лабораторной физике без гравитации сдвиг всех энергий на одну и ту же константу обычно ничего не меняет. Спектры зависят от разностей: от энергии между возбуждённым и невозбуждённым атомом или между состоянием с частицей и состоянием без неё. Кажется, что пустое пространство не даёт никакого эталона, относительно которого можно было бы измерить его абсолютную энергию.

Гравитация меняет постановку вопроса. Масса есть энергия, энергия является источником гравитации, и потому энергия вакуума гравитирует. Константа, невидимая в атомном переходе, может искривлять пространство-время.

Уравнение Эйнштейна отделяет геометрию от источников:
\begin{equation}
G_{\mu\nu}=8\pi G_N T_{\mu\nu},
\label{eq:l08-einstein}
\end{equation}
в единицах, где \(c=1\). Тензор Эйнштейна \(G_{\mu\nu}\) строится из метрики \(g_{\mu\nu}(x)\) и её производных. Тензор энергии-импульса \(T_{\mu\nu}\) содержит плотность энергии, плотность импульса, давление и напряжения, создаваемые веществом и негравитационными полями. Энергия вакуума относится к этой части источников точно так же, как и обычная материя.\editorialnote{Во вступительном устном пояснении \(8\pi/3\) кратко связывается с тензорным уравнением. Множитель \(1/3\) относится к приведённому ниже уравнению Фридмана для пространственно плоского случая, а не к уравнению~\eqref{eq:l08-einstein}.}

\begin{figure}[htbp]
\centering
\includegraphics[width=0.62\textwidth]{lecture_08_figure_02.png}
\caption{Геометрическая часть уравнения Эйнштейна, построенная из метрики пространства-времени \(g_{\mu\nu}(x)\).}
\lectureframe{00:21:58.820}
\end{figure}

\section{Сведение геометрии к масштабному фактору}

Выберем теперь простейшую однородную расширяющуюся геометрию с плоскими пространственными срезами:
\begin{equation}
ds^2=-dt^2+a(t)^2\,d\mathbf{x}^2,
\qquad
d\mathbf{x}^2=dx^2+dy^2+dz^2.
\label{eq:l08-frw-metric}
\end{equation}
Координатные метки \(\mathbf{x}\) нанесены на бесконечный резиновый лист. Они не обязаны двигаться, но их физические расстояния меняются. Две сопутствующие точки, разделенные на \(\Delta\mathbf{x}\), имеют при фиксированном \(t\) собственное расстояние
\begin{equation}
D(t)=a(t)\lvert\Delta\mathbf{x}\rvert.
\label{eq:l08-proper-distance}
\end{equation}
Итак, вся временная зависимость этой высокосимметричной геометрии заключена в одной функции — масштабном факторе \(a(t)\). Постоянное \(a\) дает плоское пространство-время Минковского после пространственного масштабирования. Зависящее от времени \(a\) может оставлять каждый пространственный срез плоским, делая при этом искривленным всё пространство-время.

\begin{classroomqa}
\audiencequestion{Для координатных меток \(x=7\) и \(x=13\) их координатное расстояние равно \(4\)?\lecturetimestamp{00:27:40}}
\lecturerresponse{Нет. Оно равно \(13-7=6\), так что их физическое расстояние равно \(6a(t)\). Это уточнение полезно: \(x\) задает сопутствующее положение, тогда как умножение на \(a(t)\) дает физическое расстояние.}
\end{classroomqa}

Для этой метрики и пространственно однородной плотности энергии одна компонента уравнения Эйнштейна принимает вид уравнения Фридмана для пространственно плоского случая,
\begin{equation}
\left(\frac{\dot a}{a}\right)^2
=\frac{8\pi G_N}{3}\rho(t).
\label{eq:l08-friedmann}
\end{equation}
Отношение
\begin{equation}
H(t)\equiv\frac{\dot a(t)}{a(t)}
\label{eq:l08-hubble-definition}
\end{equation}
есть параметр Хаббла. На данном этапе это определение, и оно может зависеть от времени.

\begin{figure}[htbp]
\centering
\includegraphics[width=0.90\textwidth]{lecture_08_figure_03.png}
\caption{Плоская расширяющаяся метрика и уравнение Фридмана после того, как на доске был восстановлен недостающий множитель \(1/a^2\).}
\lectureframe{00:30:30.000}
\end{figure}

\paragraph{Картина энергетического баланса.}
Для этой картины энергетического баланса уравнение~\eqref{eq:l08-friedmann} схематически записывается как
\begin{equation}
-\frac{3}{8\pi G_N}H^2+\rho=0.
\label{eq:l08-energy-heuristic}
\end{equation}
Энергия материи положительна, тогда как гравитационная степень свободы, представленная расширяющимся масштабным фактором, несет отрицательный кинетический член. В этой редуцированной задаче равенство нулю суммы отрицательной гравитационной энергии и положительной энергии материи служит удобной мнемонической схемой.

\editorialnote{Уравнение~\eqref{eq:l08-energy-heuristic} является гамильтоновым ограничением в сильно симметричной минисуперпространственной модели, а не координатно-независимой теоремой о том, что любая произвольная вселенная обладает глобально сохраняющейся полной энергией, равной нулю. В лекции оно используется как эвристическая интерпретация ограничения Фридмана.}

\section{Расширение с доминированием материи}

Вернемся к сопутствующему математическому ящику с координатной длиной стороны \(\Delta x=1\). Его физическая длина стороны равна \(a(t)\), так что его физический объем равен \(a(t)^3\). Если он содержит \(N\) нерелятивистских частиц массы \(m\), то
\begin{equation}
\rho_m(t)=\frac{Nm}{a(t)^3}
=\frac{\rho_0}{a(t)^3}.
\label{eq:l08-matter-density}
\end{equation}
Сами частицы не расширяются; расширяется их среднее расстояние друг от друга.

\begin{figure}[htbp]
\centering
\includegraphics[width=0.84\textwidth]{lecture_08_figure_04.png}
\caption{Метрика и уравнение Фридмана по мере того, как строится аргумент с сопутствующим ящиком.}
\lectureframe{00:36:07.940}
\end{figure}

Подстановка в уравнение~\eqref{eq:l08-friedmann} дает
\begin{equation}
\left(\frac{\dot a}{a}\right)^2
=\frac{8\pi G_N\rho_0}{3a^3}.
\end{equation}
Положим \(K=8\pi G_N\rho_0/3\) и выберем ветвь расширения:
\begin{align}
\dot a^2&=\frac{K}{a},\\
\dot a&=\sqrt{K}\,a^{-1/2},\\
a^{1/2}\,da&=\sqrt{K}\,dt.
\label{eq:l08-matter-separated}
\end{align}

\begin{classroomqa}
\audiencequestion{Как следует решать уравнение с разделенными переменными?\lecturetimestamp{00:38:15}}
\lecturerresponse{Проинтегрировать обе части. Интеграл от \(a^{1/2}\,da\) равен \(\frac{2}{3}a^{3/2}\), а не \(\frac{3}{2}a^{3/2}\); \(\frac{3}{2}\) — это новый показатель степени, а \(\frac{2}{3}\) — коэффициент.}
\end{classroomqa}

Соответственно,
\begin{equation}
\frac{2}{3}a^{3/2}
=\sqrt{K}\,(t-t_B),
\end{equation}
где \(t_B\) — постоянная интегрирования, задающая выбранное начало отсчета времени. С точностью до нормировки и сдвига по времени,
\begin{equation}
a(t)\propto (t-t_B)^{2/3}.
\label{eq:l08-matter-solution}
\end{equation}

\begin{figure}[htbp]
\centering
\includegraphics[width=0.86\textwidth]{lecture_08_figure_05.png}
\caption{Подстановка плотности материи и разделение переменных, приводящие к \(a^{1/2}da\propto dt\).}
\lectureframe{00:38:30.000}
\end{figure}

Кривая \(t^{2/3}\) растет медленнее прямой. Ее наклон уменьшается, поэтому вселенная с доминированием материи расширяется с замедлением. В ньютоновском языке, использованном в лекции, частицы сначала расходились, а их взаимное гравитационное притяжение непрерывно уменьшает скорость расхождения.

\begin{figure}[htbp]
\centering
\includegraphics[width=0.86\textwidth]{lecture_08_figure_06.png}
\caption{Сравнение на доске линейного расширения с вогнутым вниз законом для материи \(a(t)\propto t^{2/3}\).}
\lectureframe{00:40:00.000}
\end{figure}

Для ориентировки, для излучения имеем
\begin{equation}
\rho_r\propto a^{-4},
\end{equation}
а не \(a^{-3}\): разрежение объема дает три степени, а космологическое красное смещение добавляет еще одну. Соответствующий масштабный фактор тоже можно найти, но лекция намеренно оставляет это упражнение в стороне.

\section{Расширение с доминированием вакуумной энергии}

Теперь поместим в правую часть уравнения Фридмана только вакуумную энергию. Поскольку \(\rho_{\rm vac}\) постоянно,
\begin{equation}
\left(\frac{\dot a}{a}\right)^2
=\frac{8\pi G_N}{3}\rho_{\rm vac}
\equiv H^2.
\label{eq:l08-de-sitter-friedmann}
\end{equation}
Следует различать два употребления \(H\). Уравнение~\eqref{eq:l08-hubble-definition} определяет вообще говоря зависящую от времени величину \(\dot a/a\). Уравнение~\eqref{eq:l08-de-sitter-friedmann} утверждает, что во вселенной с чисто вакуумной энергией эта величина равна постоянной
\begin{equation}
H=\sqrt{\frac{8\pi G_N\rho_{\rm vac}}{3}}.
\end{equation}
Выбирая ветвь расширения, получаем
\begin{equation}
\dot a=Ha,
\end{equation}
и, следовательно,
\begin{equation}
a(t)=C e^{Ht}.
\label{eq:l08-exponential-solution}
\end{equation}
Постоянная интегрирования \(C\) фиксирует произвольную нормировку сопутствующих координат. Рост, пропорциональный текущему значению \(a\), есть экспоненциальный рост.

\begin{figure}[htbp]
\centering
\includegraphics[width=0.88\textwidth]{lecture_08_figure_07.png}
\caption{Уравнение при доминировании вакуумной энергии \((\dot a/a)^2=H^2\), его форма первого порядка \(\dot a=Ha\) и экспоненциальное решение.}
\lectureframe{00:47:45.000}
\end{figure}

Закон Хаббла для сопутствующих объектов имеет вид
\begin{equation}
v_{\rm rec}=HD.
\label{eq:l08-hubble-law}
\end{equation}
В этом чистом случае де Ситтера собственное расстояние, при котором скорость удаления равна \(c\), постоянно:
\begin{equation}
D_H=\frac{c}{H},
\label{eq:l08-hubble-radius}
\end{equation}
или \(D_H=H^{-1}\) в единицах \(c=1\), использовавшихся на доске. Этот фиксированный масштаб снова появится геометрически как радиус горизонта событий де Ситтера.\editorialnote{Радиус Хаббла не является автоматически горизонтом событий в произвольной зависящей от времени космологии. Это отождествление работает здесь потому, что \(H\) постоянно в точном пространстве де Ситтера.}

В лекции предлагается мысленный эксперимент для детектора вакуумной энергии. Разместите два очень легких пробных тела вдали от других источников, первоначально покоящихся относительно сопутствующего потока, и сделайте их взаимное притяжение пренебрежимо малым. Их расстояние подчиняется экспоненциальному расширению. Эффект слишком мал для земной установки, но космические расстояния и глубина взгляда в прошлое делают ту же физику наблюдаемой.

\section{От доминирования материи к доминированию вакуумной энергии}

Реальная космологическая модель, рассматриваемая в лекции, содержит оба члена:
\begin{equation}
H(t)^2
=\frac{8\pi G_N}{3}
\left(\frac{\rho_0}{a(t)^3}+\rho_{\rm vac}\right).
\label{eq:l08-mixed-friedmann}
\end{equation}
При малых \(a\) доминирует член материи, потому что \(a^{-3}\) велико, и решение напоминает \(t^{2/3}\). По мере роста вселенной материя разрежается, тогда как вакуумная энергия остается фиксированной. В конце концов постоянный член берет верх, и масштабный фактор загибается вверх к \(e^{Ht}\).

\begin{figure}[htbp]
\centering
\includegraphics[width=0.87\textwidth]{lecture_08_figure_08.png}
\caption{Качественный переход, изображенный в лекции: раннее замедляющееся расширение с доминированием материи, за которым следует поздний экспоненциальный рост с доминированием вакуумной энергии.}
\lectureframe{00:53:30.000}
\end{figure}

\begin{classroomqa}
\audiencequestion{Экспоненциальное решение содержит постоянную интегрирования. Откуда мы знаем, что она не чрезвычайно мала?\lecturetimestamp{00:52:41}}
\lecturerresponse{Мы ее измеряем. Смотреть дальше во Вселенную значит смотреть дальше в прошлое, поэтому наблюдения на разных расстояниях восстанавливают более ранний участок кривой расширения, а не только ее сегодняшнее значение.}
\end{classroomqa}

В лекции сообщается, что доступные тогда наблюдения поздней Вселенной согласовывались с экспоненциальной ветвью примерно на уровне 5 процентов и что эти данные выглядели убедительными, хотя и не столь точными, как прецизионное измерение вроде свойств электрона. Это следует читать как историческую оценку лекции, а не как современное утверждение о точности.

\begin{classroomqa}
\audiencequestion{Не приводит ли загибающаяся вверх кривая расширения в конечном счете к тому, что мы не сможем проследить наблюдения назад до Большого взрыва или космического микроволнового фона?\lecturetimestamp{00:55:05}}
\lecturerresponse{Этот вопрос не рассматривается как вопрос об экстраполяции графика. Существенный инвариантный вопрос — существует ли горизонт событий, и именно к этому причинному вопросу мы теперь переходим.}
\end{classroomqa}

\section{Связанные системы и далекое будущее}

Если экспоненциальное расширение продолжается, всякая компонента, которая разрежается, становится пренебрежимо малой по сравнению с \(\rho_{\rm vac}\). Это не означает, что сами атомы, солнечные системы или галактики растягиваются. Системы, удерживаемые достаточно сильно, чтобы оставаться связанными, не следуют потоку Хаббла.

\begin{classroomqa}
\audiencequestion{Начнут ли галактики, которые сегодня не участвуют в расширении, в конце концов расширяться, когда станет доминировать вакуумная энергия?\lecturetimestamp{00:57:44}}
\lecturerresponse{Связанные системы остаются связанными. Млечный Путь и Андромеда принадлежат к близкой структуре, которая останется вместе, тогда как достаточно далекие галактики, уже движущиеся вместе с расширением, продолжают удаляться и в конце концов уходят за наш горизонт.}
\end{classroomqa}

Важно различать \emph{притяжение} и \emph{связанность}. Ракета, запущенная со скоростью, превышающей скорость убегания, остается под действием притяжения Земли, однако не возвращается. Точно так же обычное гравитационное притяжение может быть наибольшей локальной силой между двумя галактиками, хотя их относительная орбита при этом не является связанной.

\begin{classroomqa}
\audiencequestion{Наша галактика гравитационно ассоциирована с группой. Преодолевает ли ускоренное расширение в конце концов это притяжение?\lecturetimestamp{00:59:17}}
\lecturerresponse{Вместе остаётся только действительно связанная часть локальной структуры. Описательного ярлыка «скопление» недостаточно; связанность означает, что относительное движение лежит ниже скорости убегания. Более далёкие галактики могут испытывать притяжение и тем не менее удаляться со скоростью, превышающей скорость убегания.}
\end{classroomqa}

Для наглядности в лекции относительное радиальное ускорение схематически представлено как сумма притягивающего члена, обратно пропорционального квадрату расстояния, и вакуумного члена, линейного по расстоянию:
\begin{equation}
\ddot r\sim-\frac{G_NM}{r^2}+H^2r.
\label{eq:l08-newtonian-vacuum}
\end{equation}
В идеализированном пределе де Ситтера коэффициенты не зависят от времени. Второй член растёт только потому, что растёт \(r\); он не возникает из-за постепенного уменьшения \(G_N\).

\editorialnote{Уравнение~\eqref{eq:l08-newtonian-vacuum} — это ньютоновская эвристика для относительного ускорения, а не дополнительный закон силы, фундаментальный для общей теории относительности. В точном пространстве де Ситтера соседние свободно падающие мировые линии имеют относительное ускорение, пропорциональное \(H^2r\).}

\begin{figure}[htbp]
\centering
\includegraphics[width=0.88\textwidth]{lecture_08_figure_09.png}
\caption{Сравнение на доске притяжения, обратно пропорционального квадрату расстояния, с эффективным вакуумным отталкиванием, линейно растущим с расстоянием.}
\lectureframe{01:04:10.000}
\end{figure}

\begin{classroomqa}
\audiencequestion{Усиливается ли сам направленный наружу эффект со временем?\lecturetimestamp{01:01:18}}
\lecturerresponse{При фиксированном расстоянии — нет. Его коэффициент постоянен. Свободно удаляющиеся объекты расходятся всё дальше, и тогда член, пропорциональный \(r\), больше, потому что расстояние между ними больше.}
\end{classroomqa}

\begin{classroomqa}
\audiencequestion{Можно ли описать суммарный эффект как постепенное уменьшение ньютоновской постоянной?\lecturetimestamp{01:02:30}}
\lecturerresponse{Нет. \(G_N\) остаётся постоянной. При заданном расстоянии оба коэффициента не меняются; с расстоянием меняется только относительная важность притяжения \(1/r^2\) и члена, пропорционального \(r\).}
\end{classroomqa}

\begin{classroomqa}
\audiencequestion{Были ли те же самые галактики гравитационно связанными в ранней Вселенной, до того как энергия вакуума стала существенной?\lecturetimestamp{01:05:03}}
\lecturerresponse{Преобладание обычного притяжения само по себе не означает связанности. Многие галактики уже удалялись примерно со скоростью убегания или быстрее. Позднее энергия вакуума определяет, замедляется ли крупномасштабное расширение к повторному коллапсу или, наоборот, ускоряется, но утверждения «какая сила больше» и «орбита связана» остаются разными.}
\end{classroomqa}

Таким образом, наблюдатель далёкого будущего может жить внутри слившейся локальной структуры, тогда как каждая несвязанная галактика станет ненаблюдаемой. Не имея исторических записей, такой наблюдатель мог бы вполне разумно заключить, что Вселенная почти пуста. Именно это причинное одиночество, а не разрушение локальных объектов, физически подводит нас к горизонту.

\section{Что означает горизонт событий}

После перерыва положим
\begin{equation}
a(t)=e^{Ht}
\end{equation}
и запишем метрику в виде
\begin{equation}
ds^2=-dt^2+e^{2Ht}d\mathbf{x}^2.
\label{eq:l08-de-sitter-metric}
\end{equation}

\begin{classroomqa}
\audiencequestion{Как называется пространство-время, описываемое этой метрикой?\lecturetimestamp{01:08:13}}
\lecturerresponse{Это пространство-время де Ситтера, обычно сокращённо называемое пространством де Ситтера. Его горизонты важны, потому что Вселенная с доминированием вакуумной энергии на поздних временах приближается к этой геометрии.}
\end{classroomqa}

Перед тем как рисовать его причинную структуру, напомним роль диаграммы Пенроуза. Бесконечные интервалы сжимаются в конечное изображение, тогда как светоподобные лучи остаются под углом \(45^\circ\). В плоском пространстве-времени массивные объекты оканчиваются в будущей времениподобной бесконечности, световые лучи — в светоподобной бесконечности, и у инерциального наблюдателя, который может ждать сколь угодно долго, нет горизонта событий, навсегда исключающего некоторую область пространства-времени.

\begin{figure}[htbp]
\centering
\includegraphics[width=0.80\textwidth]{lecture_08_figure_10.png}
\caption{Диаграмма Пенроуза для плоского пространства-времени, используемая, чтобы различить времениподобные и светоподобные конечные пункты перед введением горизонта.}
\lectureframe{01:10:35.000}
\end{figure}

Диаграмма чёрной дыры даёт нужный контраст. Внешний наблюдатель может получать сигналы только от событий, лежащих в причинном прошлом будущей мировой линии этого наблюдателя. Даже когда время наблюдения стремится к бесконечности, внутренняя область остаётся исключённой. Граница между событиями, которые могут сообщаться с этим наблюдателем, и событиями, которые никогда не могут, и есть горизонт событий.

\begin{definition}Будущий горизонт событий наблюдателя — это граница множества событий, из которых причинный сигнал, направленный в будущее, вообще может достигнуть мировой линии этого наблюдателя.
\label{def:l08-event-horizon}
\end{definition}

\begin{figure}[htbp]
\centering
\includegraphics[width=0.80\textwidth]{lecture_08_figure_11.png}
\caption{Сопоставление с чёрной дырой: обратные световые конусы для сколь угодно поздних внешних наблюдений никогда не включают область за горизонтом событий.}
\lectureframe{01:13:50.000}
\end{figure}

Следовательно, горизонт событий — это глобальная причинная граница, а не просто расстояние, на котором некоторый объект сегодня случайно выглядит тусклым. Чтобы определить его положение, нужно знать, какие сигналы могут прийти даже после сколь угодно долгого ожидания.

\section{Сжатие бесконечного будущего}

Пространственные координаты в уравнении~\eqref{eq:l08-de-sitter-metric} уже пробегают бесконечный плоский срез. Мы оставим их без изменения и сожмём только время. Введём конформное время \(T\), потребовав
\begin{equation}
dt^2=e^{2Ht}dT^2.
\label{eq:l08-conformal-definition}
\end{equation}
Взяв корень, соответствующий направлению в будущее, и перегруппировав члены, получаем
\begin{equation}
dT=e^{-Ht}dt.
\end{equation}

\begin{classroomqa}
\audiencequestion{Чему равен интеграл, включая знак?\lecturetimestamp{01:17:25}}
\lecturerresponse{Это \(-H^{-1}e^{-Ht}\). Аддитивная константа лишь сдвигает начало отсчёта \(T\), поэтому выберем её так, чтобы бесконечное будущее приходилось на \(T=0\).}
\end{classroomqa}

Итак,
\begin{equation}
T=-\frac{1}{H}e^{-Ht},
\qquad -\infty<T<0.
\label{eq:l08-conformal-time}
\end{equation}
При \(t\to-\infty\) имеем \(T\to-\infty\). Когда собственное время \(t\to+\infty\), \(T\to0^{-}\). Бесконечное будущее помещено на конечную горизонтальную линию \(T=0\); ни один наблюдатель не достигает её за конечное собственное время.

\begin{figure}[htbp]
\centering
\includegraphics[width=0.86\textwidth]{lecture_08_figure_12.png}
\caption{Интеграл для конформного времени и отображение бесконечного будущего \(t\to+\infty\) в \(T=0^{-}\).}
\lectureframe{01:19:10.000}
\end{figure}

Из уравнения~\eqref{eq:l08-conformal-time},
\begin{equation}
HT=-e^{-Ht},
\qquad
e^{2Ht}=\frac{1}{H^2T^2}.
\label{eq:l08-conformal-factor}
\end{equation}
Подстановка \(dt=e^{Ht}dT\) в уравнение~\eqref{eq:l08-de-sitter-metric} затем даёт
\begin{equation}
ds^2=\frac{1}{H^2T^2}
\left(-dT^2+d\mathbf{x}^2\right).
\label{eq:l08-conformal-metric}
\end{equation}
Выражение в скобках — это метрика Минковского. Вся геометрия де Ситтера теперь заключена в общем конформном множителе, который расходится при \(T\to0^{-}\).

\begin{figure}[htbp]
\centering
\includegraphics[width=0.87\textwidth]{lecture_08_figure_13.png}
\caption{Алгебра, связывающая \(e^{2Ht}\) с \(1/(H^2T^2)\) и подготавливающая конформный вид метрики.}
\lectureframe{01:22:50.000}
\end{figure}

\begin{classroomqa}
\audiencequestion{Какое условие определяет траекторию светового луча в этой метрике?\lecturetimestamp{01:23:49}}
\lecturerresponse{Световой луч идёт по нулевой траектории, \(ds^2=0\). Поэтому общий конформный множитель сокращается, и вдоль выбранного пространственного направления остаётся \(dT=\pm dx\).}
\end{classroomqa}

На диаграмме \((T,x)\) световые лучи являются прямыми линиями под углом \(45^\circ\):
\begin{equation}
\frac{dx}{dT}=\pm1.
\label{eq:l08-null-lines}
\end{equation}
Координатная картина сильно сжата вблизи \(T=0\), но она в точности сохраняет то, что нам нужно: какие события могут сообщаться.

\begin{figure}[htbp]
\centering
\includegraphics[width=0.86\textwidth]{lecture_08_figure_14.png}
\caption{Конформная метрика в тот момент, когда условие нулевого интервала \(ds^2=0\) используется для получения световых лучей под углом \(45^\circ\).}
\lectureframe{01:24:05.000}
\end{figure}

\section{Алиса и Боб в пространстве де Ситтера}

Поместим Алису в фиксированную сопутствующую точку \(x=0\). Её мировая линия поднимается вертикально и оканчивается на диаграмме при \(T=0\), хотя её наручным часам требуется бесконечное собственное время, чтобы туда добраться. Проведём из этой конечной точки направленные в прошлое светоподобные лучи. Они ограничивают наибольшую область, из которой Алиса вообще может когда-либо получить сигнал. События вне её остаются навсегда недоступными.

\begin{figure}[htbp]
\centering
\includegraphics[width=0.88\textwidth]{lecture_08_figure_15.png}
\caption{Предельный прошлый световой конус Алисы. Его светоподобная граница — это её зависящий от наблюдателя горизонт событий де Ситтера.}
\lectureframe{01:28:50.000}
\end{figure}

Этот горизонт индивидуален. У другого наблюдателя другая будущая конечная точка и, следовательно, другой предельный прошлый световой конус. Предположим, что в какой-то момент Боб находится внутри причинной области Алисы и они обмениваются сообщениями. Если позднее Боб, следуя потоку Хаббла, пересекает горизонт Алисы, сигналы, испущенные после этого пересечения, не могут достичь Алисы, сколько бы она ни ждала. У Боба есть собственный горизонт, и со временем он утрачивает соответствующую возможность принимать поздние сигналы от Алисы.

\begin{figure}[htbp]
\centering
\includegraphics[width=0.88\textwidth]{lecture_08_figure_16.png}
\caption{Алиса и Боб с их различными причинными границами. Сопутствующий Боб может выйти из причинного прошлого Алисы, хотя раньше они обменивались сигналами.}
\lectureframe{01:30:20.000}
\end{figure}

Если Алиса и Боб держатся за руки или, реалистичнее, остаются связанными при фиксированном собственном расстоянии, их разделение по сопутствующим координатам убывает как \(a^{-1}\). На конформной диаграмме их кривые сближаются, но это не означает, что их собственное расстояние обращается в нуль. Растущий конформный множитель компенсирует сокращение координатного промежутка, и связь остаётся возможной сколь угодно долго.

\begin{classroomqa}
\audiencequestion{Может ли Боб видеть Алису в тот же самый момент после того, как их разделят горизонты?\lecturetimestamp{01:31:43}}
\lecturerresponse{Ни один наблюдатель никогда не видит далёкое событие «в тот же самый момент» в этом буквальном смысле; наблюдение всегда относится к событию в прошлом световом конусе. Более точное утверждение состоит в том, что после пересечения горизонта никакой вновь испущенный сигнал с другой мировой линии не может прийти.}
\end{classroomqa}

В космологическом приложении к далёкому будущему Алиса представляет наш локальный связанный остаток, а Боб — каждую несвязанную далёкую галактику. Галактики не исчезают внезапно в некоторый конечный момент по часам Алисы. Их более поздние сигналы растягиваются на всё более длинные интервалы и смещаются ко всё большим длинам волн, пока не становятся физически необнаружимыми.

\begin{classroomqa}
\audiencequestion{Действительно ли теория говорит, что в далёком будущем не останется видимых далёких галактик?\lecturetimestamp{01:33:38}}
\lecturerresponse{Несвязанные галактики уходят за горизонт событий локального наблюдателя. Их последние видимые сигналы приходят со всё большей задержкой и испытывают всё большее красное смещение, так что для практических целей наблюдаемой остаётся только локальная связанная структура.}
\end{classroomqa}

\section{Красное смещение на горизонте}

На диаграмме всё ещё можно провести старый световой сигнал от Боба, достигающий Алисы при очень позднем конформном времени, из-за чего может показаться, что она продолжает его видеть. Физическое содержание заключено в интервалах между этими сигналами и в их энергии. Равные интервалы на часах Боба приходят, будучи разделены всё большими интервалами собственного времени Алисы. Частоты и энергии фотонов убывают:
\begin{equation}
\nu_{\rm received}\longrightarrow0,
\qquad
E_{\rm received}=h\nu_{\rm received}\longrightarrow0.
\end{equation}
Изображение Боба замедляется, краснеет и тускнеет. Предельный сигнал несёт исчезающе мало пригодной к использованию информации.

\begin{figure}[htbp]
\centering
\includegraphics[width=0.89\textwidth]{lecture_08_figure_17.png}
\caption{Состояние доски в конце лекции: сравнение красного смещения на горизонте де Ситтера с привычной внешней картиной горизонта чёрной дыры.}
\lectureframe{01:36:20.000}
\end{figure}

\begin{classroomqa}
\audiencequestion{Похоже ли это на чёрную дыру, где внешний наблюдатель никогда не видит, как падающий объект пересекает горизонт?\lecturetimestamp{01:36:09}}
\lecturerresponse{Да. В обоих случаях наблюдатель получает сигналы со всё большим красным смещением и никогда непосредственно не видит завершения пересечения. Причинная диаграмма показывает, какие более поздние излучённые сигналы никогда не смогут прийти; красное смещение объясняет, как затухает последняя наблюдаемая информация.}
\end{classroomqa}

Эта аналогия вызывает вопрос о температуре.

\begin{classroomqa}
\audiencequestion{Как Алиса могла бы измерить температуру непосредственно над горизонтом?\lecturetimestamp{01:39:34}}
\lecturerresponse{Операционально она могла бы спустить термометр на тросе, не дать ему пересечь горизонт и затем поднять обратно. Наблюдатель в пространстве де Ситтера мог бы представить себе аналогичный эксперимент вблизи космологического горизонта. В квантовой теории горизонты обладают температурой, но в этой лекции мы останавливаемся на классической геометрии и откладываем этот расчёт.}
\end{classroomqa}

Никакая формула для температуры де Ситтера здесь не выводится. Излучение Хокинга упоминается лишь как предварительное указание на квантовую поправку к этой классической картине.

\section{Постоянный радиус горизонта}

Конформную диаграмму можно перевести обратно в исходную временную координату. Для двух сопутствующих точек с фиксированным координатным расстоянием выполняется
\begin{equation}
D(t)=e^{Ht}\lvert\Delta x\rvert,
\end{equation}
поэтому их собственное расстояние растёт экспоненциально.

\begin{classroomqa}
\audiencequestion{Как выглядит то же движение в исходной, непреобразованной временной координате?\lecturetimestamp{01:41:14}}
\lecturerresponse{При фиксированных сопутствующих координатах вся пространственная метрика растёт как \(e^{2Ht}\); следовательно, собственные расстояния растут как \(e^{Ht}\). Если обратить преобразование координат, это расширение становится явным.}
\end{classroomqa}

Горизонт ведёт себя иначе, потому что сопутствующее координатное расстояние от Алисы до него не фиксировано. Из нулевой линии под углом \(45^\circ\), оканчивающейся в \(T=0\), следует, что расстояние до горизонта на срезе, где \(T<0\), равно
\begin{equation}
\Delta x_H=-T=\frac{1}{H}e^{-Ht}.
\label{eq:l08-coordinate-horizon}
\end{equation}
Умножая на масштабный фактор, получаем
\begin{equation}
D_H(t)=a(t)\Delta x_H
=e^{Ht}\frac{e^{-Ht}}{H}
=\frac{1}{H}.
\label{eq:l08-fixed-horizon}
\end{equation}
Восстанавливая единицы, имеем \(D_H=c/H\). Координатный зазор уменьшается, масштаб метрики растёт, а их произведение остаётся постоянным.

\begin{classroomqa}
\audiencequestion{С течением времени расстояние от Алисы до горизонта растёт, уменьшается или остаётся тем же?\lecturetimestamp{01:42:42}}
\lecturerresponse{Оно остаётся тем же. Его координатное расстояние уменьшается, тогда как конформный множитель растёт, так что собственный радиус горизонта остаётся фиксированным и равным \(H^{-1}\) в единицах, где \(c=1\).}
\end{classroomqa}

\begin{figure}[H]
\centering
\includegraphics[width=0.90\textwidth]{lecture_08_figure_18.png}
\caption{Итоговое состояние доски: у последовательных горизонтальных срезов координатное расстояние до горизонта уменьшается, но собственное расстояние остаётся тем же.}
\lectureframe{01:43:20.000}
\end{figure}

\begin{classroomqa}
\audiencequestion{Продолжает ли Алиса видеть прошлое, спроецированное на горизонт, потому что старые изображения так и не исчезают полностью?\lecturetimestamp{01:43:45}}
\lecturerresponse{Классически изображения сохраняются лишь в предельном смысле: из-за неограниченно растущего красного смещения их всё труднее и труднее обнаружить. В этом утверждении опущено излучение Хокинга, связанное с горизонтом, которое изменит квантовое описание.}
\end{classroomqa}

Тем самым лекция завершается на чётко определённой границе. Энергия вакуума породила экспоненциальное расширение; экспоненциальное расширение породило фиксированный, зависящий от наблюдателя горизонт событий; а частоты сигналов вблизи этого горизонта стали исчезающе малы. Температура горизонта и излучение Хокинга относятся к тому, что будет дальше. Предполагая, как говорит Susskind, что мы не пересекаем горизонты друг друга, встретимся снова на следующей лекции.
